\section{妥当性への脅威}

\paragraph{公開データの偏り}
RanSAP/RanSMAP は初期検証に有用だが、ransomware family、OS、volume condition、
benign workload の範囲は限定的である。headline accuracy ではなく、
family holdout、condition holdout、benign workload holdout を使う必要がある。

\paragraph{entropy 依存}
write entropy が高い ransomware は検出しやすい一方、既に暗号化されたボリューム、
圧縮済み benign file、backup、deduplication では entropy が信頼できない。
また、dataset-provided write entropy は embedded collector の cheap telemetry と
同一ではない。metadata-only variant を主結果として扱えるかが鍵になる。

\paragraph{workload intensity の混同}
AE は入力分布全体を再構成するため、ransomware ではなく単に write intensive な
作業を異常とみなす危険がある。これを防ぐには intensity-only baseline と
channel-wise attribution が必要である。

\paragraph{10 秒窓の遅延}
10 秒統計は装置負荷を抑えるが、attack start 直後の小さなシグナルを隠す可能性がある。
time to detect だけでなく、alert 前に失われる LBA fraction を測る。

\paragraph{MNN 変換差}
offline framework で良好な AE が、MNN 変換や量子化後に score ordering を保つとは
限らない。MNN parity は研究後半の付録ではなく、実装主張の中心評価項目である。
